\chapter{Boolean Analysis --- Gaussian Density}
%%% generated by the blueprint pipeline from TCSlib.BooleanAnalysis.GaussianSpace.GaussianDensity
\section{Overview}
Translated Gaussian densities represent affine Gaussian indicator averages and admit integrable bounds for their first two mean derivatives. These estimates give differentiability, global derivative bounds, and Lipschitz regularity of bounded measurable Gaussian averages.

\section{Declarations}

\begin{lemma}[Affine Gaussian indicator density formula]
\lean{BooleanAnalysis.GaussianSpace.integral_indicator_gaussian_affine}
\leanok
\uses{BooleanAnalysis.ThresholdFunctions.standardGaussianMeasure}
\difficulty{4}
Let $A\subseteq\mathbb R$ be measurable, $c,s\in\mathbb R$, and $s\ne0$. For $Z\sim N(0,1)$ and the Gaussian density $p_{s^2}(c,z)$ with mean $c$ and variance $s^2$, $\mathbb E1_A(c+sZ)=\int_{\mathbb R}p_{s^2}(c,z)1_A(z)\,dz$.
\end{lemma}

\begin{lemma}[Mean derivative of the Gaussian kernel]
\lean{BooleanAnalysis.GaussianSpace.hasDerivAt_gaussianPDFReal_mean}
\leanok
\difficulty{3}
For $v\ge0$, put $p_v(c,z)=(2\pi v)^{-1/2}\exp(-(z-c)^2/(2v))$ when $v>0$, and $p_0(c,z)=0$; take division by zero to be zero. For all real $c,z$, the derivative at $c$ of $m\mapsto p_v(m,z)$ is $((z-c)/v)p_v(c,z)$.
\end{lemma}

\begin{lemma}[Integrable envelope for first mean derivatives]
\lean{BooleanAnalysis.GaussianSpace.gaussianPDFReal_meanDerivative_dominated}
\leanok
\difficulty{6}
Let $v>0$ and $p_v(c,z)=(2\pi v)^{-1/2}\exp(-(z-c)^2/(2v))$. For every $M\ge0$, there exists a Lebesgue-integrable $b:\mathbb R\to\mathbb R$ such that $|((z-m)/v)p_v(m,z)|\le b(z)$ for all $z\in\mathbb R$ and all $m$ with $|m|\le M$.
\end{lemma}

\begin{lemma}[Differentiation of bounded Gaussian averages]
\lean{BooleanAnalysis.GaussianSpace.hasDerivAt_integral_gaussianPDFReal_mean}
\leanok
\uses{BooleanAnalysis.GaussianSpace.gaussianPDFReal_meanDerivative_dominated,BooleanAnalysis.GaussianSpace.hasDerivAt_gaussianPDFReal_mean}
\difficulty{5}
Let $v>0$ and $p_v(c,z)=(2\pi v)^{-1/2}\exp(-(z-c)^2/(2v))$. Let $f:\mathbb R\to\mathbb R$ be measurable with $|f(z)|\le1$ for all $z$. For every $c\in\mathbb R$, the derivative at $c$ of $m\mapsto\int p_v(m,z)f(z)\,dz$ is $\int((z-c)/v)p_v(c,z)f(z)\,dz$.
\end{lemma}

\begin{lemma}[Second mean derivative of the Gaussian kernel]
\lean{BooleanAnalysis.GaussianSpace.hasDerivAt_gaussianPDFReal_meanDerivative}
\leanok
\uses{BooleanAnalysis.GaussianSpace.hasDerivAt_gaussianPDFReal_mean}
\difficulty{3}
For $v\ge0$, put $p_v(c,z)=(2\pi v)^{-1/2}\exp(-(z-c)^2/(2v))$ when $v>0$, and $p_0(c,z)=0$; take division by zero to be zero. For real $c,z$, the derivative at $c$ of $m\mapsto((z-m)/v)p_v(m,z)$ is $((z-c)^2/v^2-1/v)p_v(c,z)$.
\end{lemma}

\begin{lemma}[Integrable envelope for second mean derivatives]
\lean{BooleanAnalysis.GaussianSpace.gaussianPDFReal_secondMeanDerivative_dominated}
\leanok
\difficulty{6}
Let $v>0$ and $p_v(c,z)=(2\pi v)^{-1/2}\exp(-(z-c)^2/(2v))$. For every $M\ge0$, there is a Lebesgue-integrable $b:\mathbb R\to\mathbb R$ such that $|((z-m)^2/v^2-1/v)p_v(m,z)|\le b(z)$ for all real $z$ and all $m$ with $|m|\le M$.
\end{lemma}

\begin{lemma}[Second differentiation of bounded Gaussian averages]
\lean{BooleanAnalysis.GaussianSpace.hasDerivAt_integral_gaussianPDFReal_meanDerivative}
\leanok
\uses{BooleanAnalysis.GaussianSpace.gaussianPDFReal_meanDerivative_dominated,BooleanAnalysis.GaussianSpace.gaussianPDFReal_secondMeanDerivative_dominated,BooleanAnalysis.GaussianSpace.hasDerivAt_gaussianPDFReal_meanDerivative}
\difficulty{5}
Let $v>0$ and $p_v(c,z)=(2\pi v)^{-1/2}\exp(-(z-c)^2/(2v))$. If $f:\mathbb R\to\mathbb R$ is measurable and $|f|\le1$, then for each real $c$ the derivative at $c$ of $m\mapsto\int((z-m)/v)p_v(m,z)f(z)\,dz$ is $\int((z-c)^2/v^2-1/v)p_v(c,z)f(z)\,dz$.
\end{lemma}

\begin{lemma}[Centered second Gaussian moment]
\lean{BooleanAnalysis.GaussianSpace.integral_centered_sq_gaussianReal}
\leanok
\difficulty{3}
For every $c\in\mathbb R$ and $v\ge0$, a random variable $Z\sim N(c,v)$ satisfies $\mathbb E(Z-c)^2=v$, including the point-mass law when $v=0$.
\end{lemma}

\begin{lemma}[Absolute Gaussian density integral bound]
\lean{BooleanAnalysis.GaussianSpace.abs_integral_gaussianPDFReal_mul_le}
\leanok
\difficulty{4}
Let $v>0$ and $p_v(c,z)=(2\pi v)^{-1/2}\exp(-(z-c)^2/(2v))$. Let $c\in\mathbb R$ and $g,b:\mathbb R\to\mathbb R$. If $b$ is integrable for $N(c,v)$ and $|g(z)|\le b(z)$ for every $z$, then $|\int p_v(c,z)g(z)\,dz|\le\int b(z)\,dN(c,v)(z)$, with a nonintegrable integral on the left assigned value zero.
\end{lemma}

\begin{lemma}[Uniform second mean-derivative bound]
\lean{BooleanAnalysis.GaussianSpace.abs_integral_gaussianPDFReal_secondMeanDerivative_le}
\leanok
\uses{BooleanAnalysis.GaussianSpace.integral_centered_sq_gaussianReal}
\difficulty{5}
Let $v>0$ and $p_v(c,z)=(2\pi v)^{-1/2}\exp(-(z-c)^2/(2v))$. For any function $f:\mathbb R\to\mathbb R$ with $|f|\le1$ and any real $c$, $|\int((z-c)^2/v^2-1/v)p_v(c,z)f(z)\,dz|\le2/v$, where nonintegrable integrals are assigned zero.
\end{lemma}

\begin{lemma}[Uniform first mean-derivative bound]
\lean{BooleanAnalysis.GaussianSpace.abs_integral_gaussianPDFReal_meanDerivative_le}
\leanok
\uses{BooleanAnalysis.GaussianSpace.abs_integral_gaussianPDFReal_mul_le,BooleanAnalysis.GaussianSpace.integral_centered_sq_gaussianReal}
\difficulty{5}
Let $v>0$ and $p_v(c,z)=(2\pi v)^{-1/2}\exp(-(z-c)^2/(2v))$. For any function $f:\mathbb R\to\mathbb R$ with $|f|\le1$ and any real $c$, $|\int((z-c)/v)p_v(c,z)f(z)\,dz|\le1+1/v$, where nonintegrable integrals are assigned zero.
\end{lemma}

\begin{lemma}[Lipschitz first mean derivative]
\lean{BooleanAnalysis.GaussianSpace.lipschitzWith_integral_gaussianPDFReal_meanDerivative}
\leanok
\uses{BooleanAnalysis.GaussianSpace.abs_integral_gaussianPDFReal_secondMeanDerivative_le,BooleanAnalysis.GaussianSpace.hasDerivAt_integral_gaussianPDFReal_meanDerivative}
\difficulty{3}
Let $v>0$ and $p_v(c,z)=(2\pi v)^{-1/2}\exp(-(z-c)^2/(2v))$. For measurable $f:\mathbb R\to\mathbb R$ with $|f|\le1$, the function $c\mapsto\int((z-c)/v)p_v(c,z)f(z)\,dz$ is globally $2/v$-Lipschitz.
\end{lemma}

\begin{lemma}[Lipschitz Gaussian averages in the mean]
\lean{BooleanAnalysis.GaussianSpace.lipschitzWith_integral_gaussianPDFReal_mean}
\leanok
\uses{BooleanAnalysis.GaussianSpace.abs_integral_gaussianPDFReal_meanDerivative_le,BooleanAnalysis.GaussianSpace.hasDerivAt_integral_gaussianPDFReal_mean}
\difficulty{3}
Let $v>0$ and $p_v(c,z)=(2\pi v)^{-1/2}\exp(-(z-c)^2/(2v))$. For measurable $f:\mathbb R\to\mathbb R$ with $|f|\le1$, the function $c\mapsto\int p_v(c,z)f(z)\,dz$ is globally $(1+1/v)$-Lipschitz.
\end{lemma}

\begin{lemma}[Continuous differentiability of Gaussian averages]
\lean{BooleanAnalysis.GaussianSpace.contDiff_one_integral_gaussianPDFReal_mean}
\leanok
\uses{BooleanAnalysis.GaussianSpace.hasDerivAt_integral_gaussianPDFReal_mean,BooleanAnalysis.GaussianSpace.lipschitzWith_integral_gaussianPDFReal_meanDerivative}
\difficulty{4}
Let $v>0$ and $p_v(c,z)=(2\pi v)^{-1/2}\exp(-(z-c)^2/(2v))$. For measurable $f:\mathbb R\to\mathbb R$ with $|f|\le1$, the function $c\mapsto\int p_v(c,z)f(z)\,dz$ is continuously differentiable on $\mathbb R$.
\end{lemma}
